\folie{1}{Begrüßung und Thema}{0:00--1:00}{1:00}{Nach der Worterteilung beginnen. Die Namen ruhig aussprechen, kurz zur Kommission schauen.}

\begin{sprechtext}

Guten Tag, Herr Professor Strunz, Frau Professor Kowal, Herr Professor Opferkuch und Herr Professor Li. Herzlich willkommen auch allen Gästen.

Vielen Dank an die Kommission für die Begutachtung und die Durchführung der Prüfung. Frau Professor Kowal, Ihnen danke ich besonders für die Betreuung meiner Promotion. Herrn Professor Opferkuch und Herrn Professor Li danke ich auch für die Anreise aus Nürnberg und Aachen.

Herr Professor Opferkuch, Sie haben bereits meine Bachelorarbeit zu Lithium-Ionen-Akkus an der TH Nürnberg betreut. Ich freue mich, dass Sie auch heute wieder dabei sind.

In meiner Dissertation geht es um neuronale Zustandsbestimmung für eingebettete Batteriemanagementsysteme. Ich zeige heute, wie sich Schätzgenauigkeit, Robustheit und Hardwarebedarf gemeinsam betrachten lassen.

\end{sprechtext}

\folie{2}{Warum Zustandsbestimmung direkt im BMS?}{1:00--1:45}{0:45}{Auf die lokalen Speicher und danach auf die Ressourcenbeschränkung zeigen.}

\begin{sprechtext}

Der Ausgangspunkt sind Batteriespeicher, die zuverlässig vor Ort arbeiten müssen. Das betrifft zum Beispiel dezentrale Energiesysteme, wie hier im MG-Farm-Kontext, und mobile Anwendungen außerhalb einer festen Infrastruktur.

Das Batteriemanagementsystem muss wissen, wie viel Ladung noch verfügbar ist und wie stark die Batterie bereits gealtert ist. Diese Größen lassen sich nicht einfach direkt messen. Sie müssen aus Spannung, Strom, Temperatur und der Betriebshistorie geschätzt werden.

Gleichzeitig stehen nur begrenzter Speicher und begrenzte Rechenleistung zur Verfügung. Auch eine dauerhafte Internetverbindung kann man nicht voraussetzen. Deshalb soll die Zustandsbestimmung direkt auf dem Mikrocontroller funktionieren.

\end{sprechtext}

\folie{3}{Drei Forschungsfragen}{1:45--2:35}{0:50}{Die Fragen in der Reihenfolge Modellgestaltung, Zuverlässigkeit und Umsetzung führen.}

\begin{sprechtext}

Daraus ergeben sich drei Forschungsfragen. Erstens: Wie kann ich zeitabhängige Alterung mit einer möglichst einfachen Modellstruktur abbilden? Dabei geht es besonders darum, welche Informationen ich dem Modell als Eingang zur Verfügung stelle.

Zweitens: Wie zuverlässig sind verschiedene Schätzer, wenn Messwerte gestört sind oder der Anfangszustand nicht stimmt? Ein kleiner Fehler unter idealen Bedingungen beantwortet diese Frage noch nicht.

Drittens: Wie lassen sich neuronale Modelle auf einem Mikrocontroller ausführen und dort hinsichtlich Speicher und Laufzeit verbessern?

Diese drei Fragen bilden den Aufbau meines Vortrags: zunächst die Eingabedarstellung, dann der Robustheitsvergleich und schließlich die Optimierung für die Hardware.

\end{sprechtext}

\folie{4}{Einordnung des eigenen Beitrags}{2:35--3:15}{0:40}{Je Spalte einen Beitrag nennen. Die Literaturangaben nicht vorlesen.}

\begin{sprechtext}

Für alle drei Bereiche gibt es bereits etablierte Ansätze. Der Beitrag meiner Arbeit liegt deshalb in den konkreten Untersuchungen und ihrer Verbindung.

Im ersten Teil analysiere ich systematisch die Eingangsmerkmale und die zeitliche Historie für die SOH-Schätzung mit einem MLP. Im zweiten Teil vergleiche ich unterschiedliche Schätzerklassen unter denselben Messstörungen. Im dritten Teil untersuche ich Pruning und Quantisierung für SOC- und SOH-Modelle auf derselben Mikrocontrollerplattform.

Dabei betrachte ich die Qualität der Schätzung gemeinsam mit dem Speicherbedarf und der gemessenen Ausführungszeit.

\end{sprechtext}

\folie{5}{NMC-Daten und Alterungsreferenz}{3:15--3:50}{0:35}{Einen Verlauf in der Grafik zeigen. Keine einzelnen Szenarien aufzählen.}

\begin{sprechtext}

Die erste Untersuchung basiert auf vierzehn NMC-Zellen in sieben Alterungsszenarien. Die Kurven zeigen, dass sich die Kapazität je nach Betriebsbedingungen unterschiedlich entwickelt. Einfluss haben unter anderem Temperatur, Ladezustand, Entladetiefe und Entladerate.

Der SOH wird hier über wiederholte Kapazitätsmessungen referenziert. Er beschreibt die verbleibende Kapazität im Verhältnis zur Nennkapazität. Eine gealterte Zelle kann also vollständig geladen sein und trotzdem weniger Ladung speichern als eine neue.

Diesen Datensatz verwende ich für die Entwicklung der SOH-Modelle.

\end{sprechtext}

\folie{6}{LFP-Daten für Benchmark und Optimierung}{3:50--4:25}{0:35}{Den Wechsel des Datensatzes deutlich machen.}

\begin{sprechtext}

Für den Robustheitsvergleich und die spätere Embedded-Optimierung verwende ich einen zweiten Datensatz mit fünfzehn LFP-Zellen. Hier wurden Lade- und Entladerate sowie Entladetiefe systematisch variiert, bei einer Umgebungstemperatur von fünfundzwanzig Grad.

Auch hier bilden Kapazitätsmessungen die Referenz für den Alterungszustand. Die verschiedenen Verläufe ermöglichen es, Modelle über unterschiedliche Belastungen und Alterungsstände hinweg zu untersuchen.

Wichtig für die Einordnung der folgenden Ergebnisse ist also: Der erste Modellteil verwendet NMC-Daten, die beiden weiteren Untersuchungen verwenden LFP-Daten.

\end{sprechtext}

\folie{7}{BMS-Plattform und Benchmark-Hardware}{4:25--5:05}{0:40}{Die beiden markierten Bausteine zeigen. H755 und H753 unterscheiden.}

\begin{sprechtext}

Die entwickelte BMS-Plattform verbindet Messung, Überwachung und eigene Zustandsmodelle. Der hier markierte BMS-Baustein misst Zellspannungen und Temperaturen und ermöglicht den passiven Ladungsausgleich.

Auf dem STM32H755 sind die Aufgaben auf zwei Kerne verteilt. Der Cortex-M4 übernimmt die hardwarenahen Aufgaben wie Messwerterfassung und Kommunikation. Der Cortex-M7 ist für die Zustandsmodelle vorgesehen.

Die später gezeigten Speicher- und Laufzeitmessungen wurden auf einem STM32H753 durchgeführt. Die eigene BMS-Platine zeigt damit die Integrationsumgebung, während der H753 die konkrete Plattform für den Hardwarebenchmark ist.

\end{sprechtext}

\folie{8}{Erster Untersuchungsteil: Eingabedarstellung}{5:05--5:35}{0:30}{Den Schwerpunkt auf Feature Engineering legen, nicht auf Gate-Gleichungen.}

\begin{sprechtext}

Der erste Untersuchungsteil setzt vor der Wahl einer aufwendigen Architektur an. Die Frage ist: Wie viel zeitabhängiges Batterieverhalten lässt sich bereits über eine geeignete Eingabedarstellung erfassen?

Als Ausgangspunkt dient ein MLP, also ein vollständig verbundenes neuronales Netz. Es besitzt keinen internen zeitlichen Zustand. Deshalb untersuche ich, wie abgeleitete Merkmale, Zeitauflösung und eine explizit zusammengestellte Historie die SOH-Schätzung beeinflussen.

\end{sprechtext}

\folie{9}{Wie eine Lag Sequence entsteht}{5:35--6:30}{0:55}{In der Grafik von den einzelnen Messgrößen zum gemeinsamen Eingangsvektor gehen.}

\begin{sprechtext}

Die Historie wird über sogenannte Lag Sequences bereitgestellt. Dafür fasse ich die Messwerte zunächst zu Zeitintervallen zusammen. Aus jedem Intervall entstehen Merkmale aus Spannung, Temperatur und kumuliertem Lade- beziehungsweise Entladestrom.

Mehrere aufeinanderfolgende Intervalle werden anschließend zu einem einzigen Eingangsvektor zusammengefügt. Das MLP bekommt damit sowohl aktuelle als auch vergangene Informationen. Der zugehörige Zielwert ist der SOH am Ende des Fensters.

Die zeitliche Abdeckung ergibt sich aus der Anzahl der Schritte und ihrer Dauer. Sechzehn Schritte mit jeweils zehn Minuten ergeben hier eine Historie von hundertsechzig Minuten. Die zeitliche Information steckt damit in den Eingängen; sie wird nicht durch einen rekurrenten Zustand im Netz gespeichert.

\end{sprechtext}

\folie{10}{Lernen, Modellwahl und unabhängiger Test}{6:30--7:20}{0:50}{Zelltrennung zuerst erklären, danach MSE und MAE unterscheiden.}

\begin{sprechtext}

Für Training und Validierung werden zehn NMC-Zellen verwendet. Vier weitere Zellen bleiben für den abschließenden Test reserviert. Zwanzig Prozent der Sequenzen aus den zehn Entwicklungszellen dienen der Validierung und der Auswahl der Hyperparameter.

Die Eingänge werden einheitlich auf den Bereich zwischen null und eins skaliert. Beim Training werden die Gewichte so angepasst, dass der mittlere quadratische Fehler, also der MSE, sinkt. Die Modellsuche variiert unter anderem Tiefe, Breite und zeitliche Eingabekonfiguration.

Die abschließende Evaluation erfolgt auf den vier separaten Testzellen mit MAE und MSE. Dadurch wird sichtbar, wie gut sich das Modell auf Zellen überträgt, die nicht zum eigentlichen Training gehören.

\end{sprechtext}

\folie{11}{Zeitauflösung und Historienlänge}{7:20--8:00}{0:40}{Eine passende Region der Fehlermatrix zeigen, nicht alle Kombinationen besprechen.}

\begin{sprechtext}

Diese Auswertung zeigt den Einfluss von Zeitauflösung und Sequenzlänge. Sehr grobe Intervalle glätten zwar Messwerte, können aber relevante Lade- und Entladedynamik verlieren.

Bei einer feinen Auflösung deckt dieselbe Anzahl von Schritten dagegen weniger Zeit ab. Wenn ich gleichzeitig eine lange Historie behalten möchte, wächst der Eingangsvektor.

In der untersuchten Konfiguration ergibt sich mit sechzehn Schritten und zehn Minuten pro Schritt ein geeigneter Kompromiss. Entscheidend ist also das Zusammenspiel von zeitlicher Abdeckung und Detailgrad. Mehr Historie oder feinere Daten führen nicht automatisch zu einer besseren Schätzung.

\end{sprechtext}

\folie{12}{SOH-Ergebnisse und Grenze der Übertragbarkeit}{8:00--8:50}{0:50}{Zuerst gute Verläufe, dann Zelle 9 als Gegenbeispiel zeigen.}

\begin{sprechtext}

Bei den hier gezeigten Zellen eins und sieben folgt die SOH-Schätzung der Kapazitätsreferenz recht gut. Die mittleren absoluten Fehler liegen je nach gezeigtem Verlauf ungefähr zwischen einem Viertel und einem halben Prozentpunkt.

Zelle neun zeigt die Grenze deutlicher. Hier treten systematische und lokale Abweichungen auf, mit einem MAE von rund zweieinhalb Prozentpunkten.

Die Untersuchung zeigt damit, dass eine Feedforward-Architektur mit geeigneter Historie zeitabhängige Alterung abbilden kann. Sie zeigt aber auch, dass die Übertragung auf andere Alterungsverläufe nicht selbstverständlich ist. Diese Frage nach der Zuverlässigkeit führt zum nächsten Teil: Wie verhalten sich Schätzer unter gezielt veränderten Bedingungen?

\end{sprechtext}
